\section{Connections}\label{m:sec:connections}

Four bodies of earlier work contain the idea of Sections~\ref{m:sec:suff} to~\ref{m:sec:loc} in different clothing: Longstaff Schwartz least squares Monte Carlo (LSM), the sequential papers of Carlin, Polson and Stoffer and of Johannes, Polson and Stroud, Higdon's calibration framework, and the particle wave contrast of the sixth problem.

\subsection{Simulate forward, regress backward: Longstaff Schwartz}\label{m:sec:lsm}

Longstaff and Schwartz~\cite{longstaffschwartz2001} price American options by simulating paths of the state forward, and then, moving backward through the exercise dates, regressing the realized discounted continuation payoff on a few basis functions of the current state. The regression estimates the conditional expectation $C_t(x)=E[V_{t+1}\mid X_t=x]$, and the exercise rule compares it with the immediate payoff. Two features make it the template for GBC. First, the regression conditions on the current state $X_t$ and nothing else, because the Markov property makes $X_t$ sufficient for the future: the sufficient statistic is what the regression is a function of. Second, the conditional expectation is estimated from simulations by least squares, which is exactly how the Jiang Wong summary $E[\theta\mid Z]$ is constructed from the training pairs of Algorithm~\ref{m:alg:gbc}.

\begin{proposition}[LSM is GBC for optimal stopping; GBC is LSM for the inverse problem]\label{m:prop:lsm}
Let $(X_t)$ be Markov, $V_{t+1}$ a payoff with $EV_{t+1}^2<\infty$, and $\{L_k\}_{k\le K}$ basis functions such that the $L_k(X_t)$ are square integrable and linearly independent in $L^2$. As the number of simulated paths grows, the LSM regression coefficients $\hat\beta_t=\arg\min_\beta\sum_m(V^{(m)}_{t+1}-\sum_k\beta_kL_k(X^{(m)}_t))^2$ converge almost surely to the coefficients of the $L^2$ projection of $C_t$ onto $\mathrm{span}\{L_k\}$, and $C_t(x)$ is the posterior mean summary of Theorem~\ref{m:thm:jw} for the ``parameter'' $V_{t+1}$ given the ``data'' $X_t$. Conversely, Step 2 of Algorithm~\ref{m:alg:gbc} with the mean regression loss and basis functions of the summary $S_N$ is an LSM regression in which the roles are: paths $\leftrightarrow$ prior draws $\theta^{(m)}$, state $X_t\leftrightarrow$ summary $S_N(Z^{(m)})$, payoff $V_{t+1}\leftrightarrow$ $\theta^{(m)}$, continuation value $\leftrightarrow$ posterior mean $E[\theta\mid S_N]$. In the linear Gaussian case of Proposition~\ref{m:prop:linsuff}, with the design held fixed, the posterior mean is linear in $S_N$, so a linear basis is exact, just as LSM with a linear basis is exact when the continuation value is linear in the state.
\end{proposition}

\begin{proof}
The population problem $\min_\beta E\big(V_{t+1}-\sum_k\beta_kL_k(X_t)\big)^2$ has a unique solution $\beta^*$, given by the normal equations $E[L(X_t)L(X_t)^\top]\beta^*=E[L(X_t)V_{t+1}]$ with $L=(L_1,\dots,L_K)$, and since $V_{t+1}-C_t(X_t)$ is orthogonal in $L^2$ to every square integrable function of $X_t$, the right side equals $E[L(X_t)C_t(X_t)]$, so $\sum_k\beta^*_kL_k$ is the $L^2$ projection of $C_t$ onto $\mathrm{span}\{L_k\}$. By the strong law the empirical normal equations converge almost surely to the population ones, whose matrix is invertible, so $\hat\beta_t\to\beta^*$ almost surely. That $C_t(x)=E[V_{t+1}\mid X_t=x]$ is the posterior mean of $V_{t+1}$ given $X_t$ is the definition of the summary of Theorem~\ref{m:thm:jw} with these roles, and the converse is the same computation with $(X_t,V_{t+1})$ replaced by $(S_N(Z),\theta)$ and the paths by the training pairs of Step~1 of Algorithm~\ref{m:alg:gbc}. In the linear Gaussian case with the design held fixed, the posterior mean is $V_NS_N$ by \eqref{m:eq:linquant}, a linear function of $S_N$, so it lies in the span of the coordinates of $S_N$ and the projection onto that span is exact.
\end{proof}

Cl\'ement, Lamberton and Protter~\cite{clement2002} prove that for fixed $K$ the LSM value converges almost surely as the number of paths $M\to\infty$ to the value of the projected rule, and that as $K\to\infty$ the projected value converges to the true value; Stentoft~\cite{stentoft2004} gives rates in the two period case. The same two limit structure governs the Jiang Wong summary:

\begin{theorem}[Two limit convergence of the regression summary]\label{m:thm:lsmgbc}
Let $(\theta^{(m)},S^{(m)})_{m\le M}$ be independent draws of $(\theta,S)$ from the prior predictive law of Step~1 of Algorithm~\ref{m:alg:gbc}, with $\theta$ in $\R^D$ or in a separable Hilbert space and $E\norm{\theta}^2<\infty$, let $\Law(S)$ be the law of $S$ under it and $\{L_k\}$ an orthonormal basis of $L^2(\Law(S))$, and let $\hat m_{K,M}(s)=\sum_{k\le K}\hat\beta_kL_k(s)$ be the least squares regression of $\theta^{(m)}$ on $L_1(S^{(m)}),\dots,L_K(S^{(m)})$, coordinate by coordinate. Then for fixed $K$, $\hat m_{K,M}\to m_K:=\sum_{k\le K}\langle E[\theta\mid S],L_k\rangle L_k$ almost surely in $L^2(\Law(S))$ as $M\to\infty$, and $m_K\to E[\theta\mid S]$ in $L^2(\Law(S))$ as $K\to\infty$. If moreover $\sum_{k\le K}L_k(s)^2\le bK$ for all $s$ and $K$ and some constant $b$, as for a basis bounded by $\sqrt b$, and $K_M\to\infty$ with $K_M\log K_M/M\to0$, then $\hat m_{K_M,M}\to E[\theta\mid S]$ in $L^2(\Law(S))$ in probability. What the regression recovers is $E[\theta\mid S]$, the posterior mean given the summary, and not the posterior mean $m(Z)=E[\theta\mid Z^{(N)}]$ itself. By Theorem~\ref{m:thm:jw}(i) one always has $E[\theta\mid S]=E[m(Z)\mid S]$, and this equals $m(Z)$ if and only if $m(Z)$ is $\sigma(S)$ measurable, which is the equality case there. That extra condition is a property of the summary and not of the regression:
\begin{enumerate}[label=(\alph*),nosep]
\item if $S=m(Z)$ then $E[\theta\mid S]=S=m(Z)$ exactly, for every $N$, by Theorem~\ref{m:thm:jw}(ii);
\item if, in the linear model, $S$ is the adjoint statistic of Proposition~\ref{m:prop:linsuff} with the design held at its observed value in Step~1, or the pair of that statistic and the Gram matrix $\Phi^\top\Phi$ when the design is drawn afresh, then $S$ is sufficient and $E[\theta\mid S]=m(Z)$ exactly, for every $N$;
\item if $S=\hat S^{J_N}_N$ is the one step statistic of Definition~\ref{m:def:summary}, then Proposition~\ref{m:prop:suff} and Theorem~\ref{m:thm:jwbvm}, the latter under Assumption~\ref{m:ass:U}, give only that $\sqrt N\,\norm{E[\theta\mid S]-m(Z)}_{H^{-k-2}}\to0$ in probability.
\end{enumerate}
Under (a) or (b) the regression summary $\hat m_{K_M,M}(S)$, the center of the GBC posterior built on it, converges to the exact posterior mean $m(Z)$ in the double limit, in $L^2(\Law(S))$, that is, on average over the prior predictive law of the data and not at a fixed observed value. Under (c), which is the case of interest for a nonlinear forward map, it has it only to $o(N^{-1/2})$ in $H^{-k-2}$, which is what (G1) of Assumption~\ref{m:ass:G} asks for and all that is claimed below.
\end{theorem}

\begin{proof}
Write $L=(L_1,\dots,L_K)$, $\hat G=M^{-1}\sum_mL(S^{(m)})L(S^{(m)})^\top$ and $\beta_k=E[\theta L_k(S)]$, the coefficients of $m_K$; when $\hat G$ is invertible, $\hat\beta=\hat G^{-1}M^{-1}\sum_mL(S^{(m)})\theta^{(m)}$, coordinate by coordinate. The first limit is the strong law: $\hat G\to E[LL^\top]=\Id$ and $M^{-1}\sum_mL(S^{(m)})\theta^{(m)}\to\beta$ almost surely. The second is the $L^2$ expansion of $E[\theta\mid S]$ in the basis. For the joint limit write $\theta^{(m)}=m_K(S^{(m)})+e^{(m)}$, so that $E[L(S)e]=0$, $E\norm e^2\le E\norm\theta^2$ and $\hat\beta-\beta=\hat G^{-1}\xi$ with $\xi=M^{-1}\sum_mL(S^{(m)})e^{(m)}$; the summands of $\xi$ are independent and centered, so $E\norm\xi^2=M^{-1}E\big[\norm{L(S)}^2\norm e^2\big]\le bK\,E\norm\theta^2/M$. The matrices $M^{-1}\big(L(S^{(m)})L(S^{(m)})^\top-\Id\big)$ are independent and centered, with norm at most $bK/M$ and with the sum of their expected squares at most $(bK/M)\Id$, so the matrix Bernstein inequality, Theorem~6.1.1 of Tropp~\cite{tropp2015}, gives $E\norm{\hat G-\Id}\le\sqrt{2bK\log(2K)/M}+bK\log(2K)/(3M)$, which tends to zero along $K=K_M$. On the event $\{\norm{\hat G-\Id}\le\frac12\}$, whose probability therefore tends to one, $\norm{\hat G^{-1}}\le2$ and, by orthonormality, $\norm{\hat m_{K,M}-m_K}^2_{L^2(\Law(S))}=\norm{\hat\beta-\beta}^2\le4\norm\xi^2=O_P(K/M)$. Hence $\hat m_{K_M,M}-m_{K_M}\to0$ in probability, and $m_{K_M}\to E[\theta\mid S]$ by the second limit.

For the identification of the limit, $E[\theta\mid S]=E[m(Z)\mid S]$ is the tower property, and the measurability criterion is the equality case of Theorem~\ref{m:thm:jw}(i), read off from the decomposition $\Var(\theta\mid S)=E[V(Z)\mid S]+\Var(m(Z)\mid S)$. Case (a) is Theorem~\ref{m:thm:jw}(ii), and case (b) is Proposition~\ref{m:prop:linsuff}, whose posterior mean $V_NS_N$ is a function of $S_N$ and $\Phi^\top\Phi$. In case (c), Theorem~\ref{m:thm:jwbvm}, which needs Assumption~\ref{m:ass:U}, applied with $S=\hat S^{J_N}_N$ gives $\Wone_{1,H^{-k-2}}(\tau_N,\tau^{\rm S}_N)\to0$ in $\mathbf{P}$ probability, where $\tau_N$ and $\tau^{\rm S}_N$ are the laws of $\sqrt N(\theta-S)$ under $\Pi(\cdot\mid Z^{(N)})$ and under $\Pi(\cdot\mid S)$; their means differ by $\sqrt N(m(Z)-E[\theta\mid S])$, and $\norm{EX-EY}_B\le\Wone_{1,B}(\Law X,\Law Y)$ on a separable Banach space $B$, which gives the stated $\sqrt N$ equivalence.
\end{proof}

Theorem~\ref{m:thm:lsmgbc} explains why the choice of basis is the whole game: in LSM the basis is Laguerre or Hermite polynomials of the state, and in GBC for a linear PDE it is the identity on the adjoint statistic. The one step statistic of Section~\ref{m:sec:suff} is the choice of ``state'' that makes the regression asymptotically linear.

The American option value solves a backward parabolic free boundary problem, and LSM is a Monte Carlo solver for that backward equation: forward paths, backward regression. Our sufficient statistic for the initial condition is the adjoint equation, also backward and also parabolic, driven by the data.

\subsection{Nonlinear diffusions with jumps}\label{m:sec:jps}

The initial condition of a PDE is a static parameter of a dynamical system observed with noise, which is the setting of our sequential work on state space models and continuous time finance.

\subsubsection{Carlin, Polson and Stoffer (1992)}

Carlin, Polson and Stoffer~\cite{cps1992} introduced the Gibbs sampler for state space models with nonlinear transition and observation functions, $x_t=g_t(x_{t-1})+w_t$, $y_t=f_t(x_t)+v_t$, and non Gaussian errors represented as scale mixtures of normals with latent scales $\lambda_t$. The full conditional of $x_t$ given its neighbors, the data and the scales is proportional to the product of three densities, $p(x_t\mid x_{t-1})\,p(x_{t+1}\mid x_t)\,p(y_t\mid x_t)$, which is low dimensional and can be sampled by rejection or Metropolis when $g_t$ is nonlinear; conditional on the scales the parameter updates are Gaussian. The output is the joint smoothing distribution. For the PDE initial condition problem with model error, $u_{k+1}=u_k+\Delta tF(u_k)+\sqrt{\Delta t}\,G\xi_k$, this is the smoother for the problem of this part: the forward edge $p(u_{k+1}\mid u_k)$ is nonlinear in $u_k$ through $F$ and the backward edge is Gaussian in $u_k$, and the single move update samples $u_k$ from their product. As $G\to0$ the chain of transitions collapses to the deterministic solution map and the smoother degenerates to a sampler for $\theta$ alone, which is the Konen Nickl setting. The scale mixture device is the same one that produces the normal inverse Gaussian (NIG) robust likelihoods of Polson and Sokolov~\cite{polsonsokolov2026ns}.

\subsubsection{Johannes, Polson and Stroud (2009)}

Johannes, Polson and Stroud (JPS)~\cite{jps2009} filter latent states in continuous time jump diffusions
\begin{align}
dS_t&=\mu_s(S_t,X_t)dt+\sigma_s(S_{t-},X_{t-})dW^s_t+d\Big(\sum_{j\le N^s_t}f^s(\xi^s_j,S_{\tau_j-},X_{\tau_j-})\Big),\label{m:eq:jd1}\\
dX_t&=\mu_x(S_t,X_t)dt+\sigma_x(S_{t-},X_{t-})dW^x_t+d\Big(\sum_{j\le N^x_t}f^x(\xi^x_j,S_{\tau_j-},X_{\tau_j-})\Big),\label{m:eq:jd2}
\end{align}
from discrete observations of $S$. In this subsection $S$ is the observed price, $N^s_t$ and $N^x_t$ count jumps, and $M$ is the number of Euler substeps between observations, so that in the comparison with Theorem~\ref{m:thm:lsmgbc} drawn below, this $M$ and the particle count $P$ correspond to $K$ and $M$ there; the adjoint statistic of Proposition~\ref{m:prop:storvik} is written $S_n$.

\subsubsection{High frequency augmentation} The continuous time transition density is unavailable, so $M$ Euler substeps are simulated between observation times (``filling in the missing data'', after Pedersen~\cite{pedersen1995}, Elerian, Chib and Shephard~\cite{elerian2001} and Eraker~\cite{eraker2001}), the augmented state $X^M_{t+1}=(X_t,X_{t+1/M},\dots,X_{t+(M-1)/M})$ is propagated by an auxiliary particle filter, and the observation likelihood is Gaussian conditional on the augmented path. The discretization error of the filter is then controlled by $M$, and the convergence result of Del Moral, Jacod and Protter~\cite{delmoral2001} requires $M$ to grow slower than the particle count, for instance $M=\sqrt P$. This is the same two index limit as Theorem~\ref{m:thm:lsmgbc} and Theorem~\ref{m:thm:qnnA}, with the roles of discretization and Monte Carlo made explicit: the PDE solver's time step in Step 1 of Algorithm~\ref{m:alg:gbc} is the analogue of $1/M$, and its error enters (G1) alongside the generator's.

\subsubsection{The continuous record remark} With continuous observation of $S$, the jump times and sizes are visible as discontinuities and the quadratic variation $\frac{d}{dt}\langle S\rangle_t=\sigma_s^2(S_t,X_t)$ reveals $X_t$ whenever $\sigma_s$ is invertible in $X$, so that, in their words, there is nothing to filter. This is a sufficiency statement: in the continuous record limit the quadratic variation is a sufficient statistic for the volatility state. Section~\ref{m:sec:pr} discusses this observation in the setting of the approximate Bayes factors of Polson and Roberts~\cite{polsonroberts1994} for discretely observed diffusions, which use the discrete data alone. Its PDE analogue is the Lipschitz stability estimate at $t=0$ of Konen and Nickl~\cite{konennickl2026} and of Nickl~\cite{nicklbvm2024}: with observations dense near the initial time the initial condition is identified with a quantitative stability constant, and information is lost exponentially only through a late start of the observations (Table~\ref{m:tab:iota}). The quadratic variation statistic is to the diffusion coefficient what the adjoint statistic $w(0)$ is to the initial condition: the functional of the observed path from which the latent quantity can be read off.

\subsubsection{Sequential parameter learning requires a sufficient statistic} For parameter learning JPS use Storvik's algorithm and state the requirement plainly: it applies when the parameter posterior given the observations and latent variables up to time $t$ is a function of a low dimensional sufficient statistic $T_t=T(S_{0:t},X^M_t)$, updated recursively, and when this fails the computational gain is lost. This is the requirement of Sections~\ref{m:sec:suff} to~\ref{m:sec:jw} for the initial condition, and the adjoint statistic satisfies it in the linear case with a recursive update:

\begin{proposition}[The adjoint statistic is a Storvik statistic]\label{m:prop:storvik}
In the linear model \eqref{m:eq:design}, order the observations by time. The sufficient statistic $S_n=\Phi_{1:n}^\top Y_{1:n}$ and Gram matrix $G_n=\Phi_{1:n}^\top\Phi_{1:n}$ satisfy the recursions
\begin{align*}
S_{n+1}&=S_n+\big(s_j(t_{n+1})e_j(x_{n+1})Y_{n+1}\big)_{j\le D},\\
G_{n+1}&=G_n+\phi_{n+1}\phi_{n+1}^\top,\qquad\phi_{n+1}=\big(s_j(t_{n+1})e_j(x_{n+1})\big)_{j\le D},
\end{align*}
and the posterior after $n$ observations is $N\big((G_n+\Sigma_0^{-1})^{-1}S_n,(G_n+\Sigma_0^{-1})^{-1}\big)$. Consequently the sequential scheme that carries $(S_n,G_n)$ and, at each step, draws particles independently from $N\big((G_n+\Sigma_0^{-1})^{-1}S_n,(G_n+\Sigma_0^{-1})^{-1}\big)$ is exact: $(S_n,G_n)$ is a deterministic function of the first $n$ observations, each new observation is absorbed by the rank one update above, and every particle is an exact draw from the posterior after $n$ observations.
\end{proposition}

\begin{proof}
Additivity of $\Phi^\top Y$ and $\Phi^\top\Phi$ over observations gives the two recursions, and \eqref{m:eq:linpost} applied to the first $n$ observations identifies the posterior. For the last assertion, the recursions show that $(S_n,G_n)$ is $\sigma(Z_1,\dots,Z_n)$ measurable, and the identification of the posterior makes each particle an exact draw from it; the particles are independent and carry no weights.
\end{proof}

\begin{remark}[Why degeneracy does not arise here, and where it returns]\label{m:rem:degeneracy}
In the linear model there is no latent state, so the particle systems of Storvik~\cite{storvik2002} and Carvalho and coauthors~\cite{cjlp2010} collapse: the sufficient statistic is a deterministic function of the data rather than of a sampled state path, all weights are equal, and there is nothing to resample. With model error, that is, in the state space formulation of Carlin, Polson and Stoffer~\cite{cps1992} with $u_{k+1}=u_k+\Delta tF(u_k)+\sqrt{\Delta t}\,G\xi_k$ and $G\ne0$, the statistic depends on the sampled path, the weights are no longer equal, and the clustering of particles that Johannes, Polson and Stroud~\cite{jps2009} report for Storvik's algorithm on S\&P 500 data after the 1987 crash can occur. We do not claim, and do not prove, that particle learning avoids degeneracy outside the case of deterministic dynamics.
\end{remark}

For nonlinear forward maps the statistic of Definition~\ref{m:def:summary} is not additive, because the adjoint is linearized at $\bar\theta_N$ and the residuals depend on it. The practical filter of Polson, Stroud and M\"uller~\cite{psm2008} and the particle learning construction of Carvalho and coauthors~\cite{cjlp2010} handle this by relinearizing at checkpoints, and Proposition~\ref{m:prop:suff} says that after enough data the relinearization error is $o(N^{-1/2})$, so the statistic becomes asymptotically additive. JPS report that Storvik's algorithm works on simulated data and degenerates on S\&P 500 data after the 1987 crash, which they attribute to misspecification; in the PDE setting the analogue is a forward model that is wrong, and the discrepancy discussion of Section~\ref{m:sec:higdon} applies.

\subsubsection{Jumps} The SVJ and SVCJ models of JPS, stochastic volatility with jumps in returns and with contemporaneous jumps in returns and volatility, separate jumps from diffusive volatility by carrying the jump indicator and size in the particle and weighting by a two component mixture likelihood. The stochastic Navier Stokes model of Polson and Sokolov~\cite{polsonsokolov2026ns} with NIG innovations is the PDE version: the latent inverse Gaussian scale plays the role of the jump size, and conditional on it the update is Gaussian. What JPS show is that the filter can identify large jumps but not those within the diffusive range, and the same limitation will hold for identifying intermittent forcing in a turbulent flow from noisy observations: the sufficient statistic for the scales is the innovation, and innovations inside the Gaussian range carry no information about them.

\subsection{GBC as inverse emulation: Higdon's calibration}\label{m:sec:higdon}

The Bayesian calibration framework of Kennedy and O'Hagan~\cite{kennedyohagan2001}, as Higdon and collaborators made it operational for expensive simulators~\cite{higdon2004,higdon2008,williams2006}, is the closest methodological ancestor of Algorithm~\ref{m:alg:gbc}.

\subsubsection{Forward versus inverse emulation} Higdon's construction runs the simulator at a design $\{\theta^{(m)}\}$, fits a Gaussian process emulator of $\theta\mapsto\Gcal(\theta)$, adds a discrepancy process $\delta$, and runs Markov chain Monte Carlo (MCMC) on $(\theta,\delta)$. GBC uses the same simulations but fits the inverse map $S_N(Z)\mapsto\theta$, so the MCMC step disappears. The emulator must be accurate uniformly over the region the MCMC visits, which limits emulator based calibration to low dimensional $\theta$; the inverse map must be accurate only near the observed summary, and Remark~\ref{m:rem:affine} says it is locally affine there. There is no transfer theorem like Theorem~\ref{m:thm:main} for emulator based calibration, because emulator error enters the likelihood nonlinearly; the closest results are the emulator error bounds of Stuart and Teckentrup~\cite{stuartteckentrup2018}.

\subsubsection{High dimensional output and the summary} Higdon and coauthors~\cite{higdon2008} project image or time series output onto a principal component basis of the simulation runs. The summary \eqref{m:eq:summary} is the data side analogue: a projection onto the first $J_N$ eigenmodes of the prior, weighted by the adjoint of the linearized forward map.

\subsubsection{Process convolutions} Higdon's process convolution priors~\cite{higdon1998,higdon2002,higdon1999}, $\theta(x)=\int k(x-s)w(s)ds$ with $w$ on a coarse lattice, are physical space versions of the series \eqref{m:eq:prior} truncated to finitely many modes; the Nickl Wang constraint $D\lesssim N^{d/(2\gamma+d)}$ is a rule for how many knots a sample supports, and nonstationary kernels are a way to make the ellipsoidal localization norm of the Nickl Wang surrogate spatially adaptive.

\subsubsection{The Darcy problem in 2002} Lee and coauthors~\cite{lee2002} inferred permeability fields with Markov random field priors and MCMC, and Higdon, Lee and Holloman~\cite{higdon2003} developed coarse solver guided MCMC, the precursor of delayed acceptance and multilevel schemes~\cite{christenfox2005,dodwell2015}; this is the elliptic problem of Giordano and Nickl~\cite{giordanonickl2020} and the cousin of the Nickl Wang steady state Schr\"odinger example. The coarse fine hierarchy is a warm start for the localized training of Section~\ref{m:sec:loc}: the coarse solve supplies the center $\bar\theta_N$ about which the fine one is trained.

\subsubsection{Discrepancy and the no model error assumption} Every theorem above assumes deterministic dynamics without model error. Higdon's framework carries $\delta$, and the identifiability problem between $\theta$ and $\delta$~\cite{brynjarsdottirohagan2014,tuowu2016} is Condition 2E of Nickl in applied form: when the discrepancy can absorb a perturbation of $\theta$, the information operator is not injective. The natural formulation is hierarchical, with $\delta$ a Gaussian process of amplitude $\sigma_N\to0$; for GBC it is trivial to include $\delta$ in Step 1, and for the theory the admissible rate of $\sigma_N$ is open.

\subsubsection{Auxiliary variables} Higdon's auxiliary variable and partial decoupling methods~\cite{higdon1998b,besag1995} are the Swendsen Wang lineage (Chapter~\ref{ch:spins}) on which the slice sampling optimizer of Birge and Polson~\cite{birgepolson2012} draws, and the latent scale augmentation reconciling heavy tailed likelihoods with the Nickl Wang Gaussian structure is the same device.

\subsection{Particles, waves, and the sixth problem}\label{m:sec:hilbert}

Deng, Hani and Ma~\cite{denghanima2025} derive the Navier Stokes Fourier equations from hard sphere dynamics through the Boltzmann equation (Chapter~\ref{ch:kinetic}), and Deng and Hani~\cite{denghani2023} derive the wave kinetic equation from the cubic nonlinear Schr\"odinger equation. In both cases the continuum law is the limit of a microscopic simulation, and in both cases the limit is a statement about a typical trajectory, not about inference from observations of one trajectory. GBC is the inferential counterpart. Step 1 of Algorithm~\ref{m:alg:gbc} may simulate the microscopic system (hard spheres, or the interacting particle system whose mean field limit is the McKean Vlasov equation of Nickl, Pavliotis and Ray~\cite{nickl2025}, or the vortex system whose limit is 2D Navier Stokes) rather than the continuum PDE, and learn the map from continuum observations to the continuum initial condition or interaction potential. The propagation of chaos error~\cite{sznitman1991} then enters (G1) and must be $o(N^{-1/2})$, which requires the particle number to grow polynomially in $N$; Castre and Nickl~\cite{castrenickl2026} have verified gradient stability for the McKean Vlasov forward map, so the Nickl Wang computational theory applies there once a feasible initializer is assumed (Section~\ref{sec:gradstab}), and a Bernstein von Mises theorem for the interaction potential, if proved, would transfer by Theorem~\ref{m:thm:main}.

The dissipative dispersive contrast of Section~\ref{m:sec:nondiss} has a sixth problem reading. On the particle side, dissipation in the continuum limit comes from collisions and the Boltzmann H theorem, and it is what makes the Navier Stokes information operator a compact perturbation of the Stokes one, with the consequences of Sections~\ref{m:sec:heat} to~\ref{m:sec:ns}. On the wave side there is no H theorem for the nonlinear Schr\"odinger equation itself; the wave kinetic equation, which does have entropy production, arises only in the limit of many four wave resonant interactions, and those same interactions are what couple Fourier modes in the information operator of Section~\ref{m:sec:nondiss} once the background is not a single plane wave. For a general background, invertibility of the information operator depends on the same resonant geometry that controls whether the kinetic limit exists: the Bayesian sixth problem for waves inherits the structure of the deterministic one.

\section{What remains}\label{m:sec:open}

Besides Open Question~\ref{m:op:planewave}, five questions are left open, and they are not of the same kind. Two are hypotheses this part uses and does not prove, so the theorems that invoke them are conditional: the uniform Bernstein von Mises statement of Assumption~\ref{m:ass:U}, without which the one step statistic is asymptotically sufficient only at a fixed truth, and the Bernstein von Mises theorem under a network prior, without which Theorem~\ref{m:thm:PR} delivers contraction rates and not the Gaussian limit. Two ask for quantities the asymptotics do not supply: the finite $N$ correction, and the error of the summary obtained by iterating the Gauss Newton step. The last, model error, changes the problem rather than refining it, and it is the one on which the manifold picture is itself at stake, since a trajectory generated with model error does not lie on the solution manifold.

\begin{openquestion}[Bernstein von Mises under a network prior]
Extend Wang and Ro\v{c}kov\'a~\cite{wangrockova2020} to the inverse problem: prove Theorem~\ref{m:thm:nickl2} with the Gaussian prior \eqref{m:eq:prior} replaced by the Polson Ro\v{c}kov\'a prior, uniformly over the test functions $\psi$ of the functional limit in a Sobolev ball, by absorbing the perturbation direction $p_N(\bar\psi_{\theta_0})/\sqrt N$, $\bar\psi_{\theta_0}=T(\I^*_{\theta_0}\I_{\theta_0})^{-1}\psi$, into the last layer. With Theorem~\ref{m:thm:PR} this would give GBC under an adaptive prior all the conclusions of Theorem~\ref{m:thm:main}.
\end{openquestion}

\begin{openquestion}[Uniformity]\label{m:op:unif}
Prove Assumption~\ref{m:ass:U}, with the rate $\bar\epsilon_N\to0$ used in Theorem~\ref{m:thm:qnnA}, by tracking the constants of Theorem~2 of Nickl~\cite{nicklbvm2024} and Theorem~2.8 of Konen and Nickl~\cite{konennickl2026} over the localization sets, and do the same for the uniform second moment bound~\eqref{m:eq:momentbound} of Proposition~\ref{m:prop:abcopt}. Theorem~\ref{m:thm:u2ploc} proves Assumption~\ref{m:ass:u2p} for localized training with the summary \eqref{m:eq:trainsummary}, and Proposition~\ref{m:prop:locres} an analogue for the centered and rescaled residuals of Algorithm~\ref{m:alg:loc} when the observed design is reused; prove it for fresh designs, and compare $N(S(Y),\Sigma_N)$ with $\PiN$ for a nonlinear forward map (Remark~\ref{m:rem:u2pres}). The bound on the error of the summary in $L^2$ and in $H^{\bar\gamma}$ that Theorem~\ref{m:thm:qnnA} assumes, which Propositions~\ref{m:prop:sumrate} and~\ref{m:prop:sumrd} prove in the cases listed after that theorem, is to be proved, from the structure of the equation, for the one step summary of a nonlinear forward map other than reaction diffusion when the preliminary estimates have only the rates of Definition~\ref{m:def:summary} or \eqref{m:eq:infobound} is not known; Remark~\ref{m:rem:sumopen} says what the route of Proposition~\ref{m:prop:sumrate} needs. Another hypothesis of the same kind is the Navier Stokes regularity $\alpha>12$ of Konen and Nickl, which Schauder theory is expected to reduce to $\alpha\ge3$; the Taylor Green example, explicit up to a sparse coupling, is the natural test case. The condition of Lemma~\ref{m:lem:gapred}, under which Proposition~\ref{m:prop:gapnl} applies to a nonlinear forward map, is of the same kind: prove it for the first $J'$ coordinates of the one step summary, with a Gaussian model whose posterior covariance is asymptotically the inverse information at the prior mean, or characterize the summaries for which Proposition~\ref{m:prop:gapnl} does not apply, such as the first coordinate of the adjoint statistic at the prior mean in Proposition~\ref{m:prop:gaptoy}. For a nonlinear forward map outside the class of Proposition~\ref{m:prop:prrd}(ii) and a fixed number $J$ of modes, determine the expected information gain $\mathbf E\,I_N$ up to $o(1)$ (Remark~\ref{m:rem:prscope} names a candidate); extend Corollary~\ref{m:cor:prreal} to $J=J_N\to\infty$. For $\beta>2$, decide whether the exponent of Theorem~\ref{m:thm:qnnB} is attained over the whole class $\mathcal Q_\beta(J)$, as it is when the densities are bounded (Proposition~\ref{m:prop:qnnUB3}).
\end{openquestion}

\begin{openquestion}[Relinearized summary]
Confirm Remark~\ref{m:rem:relin}: for the summary obtained by iterating the Gauss Newton step, with the Jacobian recomputed from the data alone at each iterate, show that the $c^4$ term disappears from \eqref{m:eq:decomp} for the converged iterate, find how fast it decays with the number of steps, and determine whether, in a nonlinear problem, the resulting error falls below the error of the Laplace approximation, of order $c_3\sqrt{J^3/N}$ in the classical bounds, with $c_3$ a bound on the third derivatives of the log likelihood.
\end{openquestion}

\begin{openquestion}[Finite $N$]
Theorem~\ref{m:thm:main} is a $\sqrt N$ statement. The heat equation example shows the Gaussian formulas are accurate at $N=2000$ for a linear problem with a calibrated prior (Table~\ref{m:tab:diag}), and Section~\ref{m:sec:expns} measures the nonlinear case at a few design points in a six mode truncation. Neither is a quantification. Quantify the non Gaussian correction for the Taylor Green problem as a function of $A/\nu$ and $N$, and bound $\Wone_1(\Pihat,\PiN)$ nonasymptotically in terms of the generator approximation error for the true conditional quantile map. Section~\ref{m:sec:expns}'s failure to find the crossover in any regime tested says the answer cannot be assumed, and suggests the obstruction: the Laplace error and the GBC bias are governed by the same curvature.
\end{openquestion}

\begin{openquestion}[Model error]
Combine Theorem~\ref{m:thm:main} with a discrepancy process of amplitude $\sigma_N$ and identify the rate at which $\sigma_N$ must tend to zero for the deterministic limit to persist, connecting to the Gibbs and particle smoothers of Section~\ref{m:sec:connections}.
\end{openquestion}